% !TEX root = ../matan.tex

\section{Степенные ряды, радиус сходимости, формула Коши--Адамара.}

\begin{Definition}{Степенной ряд и радиус сходимости}{}
	\textbf{Степенным рядом} с центром в точке $z_0 \in \C$ называется ряд
	\[
	\sum_{n=0}^{\infty} a_n (z - z_0)^n, \qquad a_n \in \C.
	\]
	Число $R \in [0, +\infty]$ называется \textbf{радиусом сходимости}, если ряд сходится при $|z - z_0| < R$ и расходится при $|z - z_0| > R$. Множество $\{|z - z_0| < R\}$~--- \textbf{круг сходимости}.
\end{Definition}

\begin{Theorem}[required]{Формула Коши--Адамара}{}
	Радиус сходимости степенного ряда $\sum_{n=0}^{\infty} a_n (z - z_0)^n$ равен
	\[
	R = \frac{1}{\displaystyle\limsup_{n \to \infty} |a_n|^{1/n}},
	\]
	с соглашениями $1/0 = +\infty$ и $1/(+\infty) = 0$.
\end{Theorem}

\begin{proof}
	Без ограничения общности $z_0 = 0$. Обозначим $b_n = |a_n|^{1/n}$; ключевое тождество:
	\[
	|a_n z^n| = \bigl( |a_n|^{1/n}\,|z| \bigr)^n = (b_n |z|)^n.
	\]
	Рассмотрим три случая в зависимости от $L = \limsup_{n} b_n$.

	\textbf{Случай 1: $\{b_n\}$ не ограничена ($L = +\infty$).} Для любого $z \neq 0$ найдётся подпоследовательность $\{n_k\}$ с $b_{n_k} > 1/|z|$, откуда $(b_{n_k} |z|)^{n_k} > 1$. Общий член ряда не стремится к нулю, ряд расходится при всех $z \neq 0$, то есть $R = 0 = 1/L$.

	\textbf{Случай 2: $0 < L < +\infty$.} Полагаем $R = 1/L$.

	\emph{Сходимость при $|z| < R$.} Из $|z| < 1/L$ следует $L < 1/|z|$, поэтому найдётся $\eps > 0$ такое, что $(L + \eps)|z| < 1$. По определению верхнего предела начиная с некоторого $N_0$ выполнено $b_n \le L + \eps$, и тогда
	\[
	b_n |z| \le (L + \eps)|z| =: q < 1, \qquad \text{откуда } |a_n z^n| = (b_n |z|)^n \le q^n.
	\]
	Ряд мажорируется сходящейся геометрической прогрессией $\sum q^n$, значит сходится абсолютно.

	\emph{Расходимость при $|z| > R$.} Из $|z| > 1/L$ следует $L|z| > 1$, поэтому существует подпоследовательность $\{n_k\}$ с $b_{n_k}|z| > 1$, то есть $(b_{n_k}|z|)^{n_k} > 1$. Общий член не стремится к нулю, ряд расходится.

	\textbf{Случай 3: $L = 0$.} Для любого $z$ найдётся $N_0$ такое, что при $n \ge N_0$ выполнено $b_n |z| < 1/2$, откуда $|a_n z^n| \le (1/2)^n$. Ряд сходится абсолютно при всех $z$, то есть $R = +\infty = 1/L$.

	Во всех случаях $R = 1/L$, что и требовалось.
\end{proof}

\begin{Remark}{Поведение на границе}{}
	Формула Коши--Адамара ничего не утверждает о точках окружности $|z - z_0| = R$: там возможна сходимость во всех точках, расходимость во всех или смешанное поведение.
\end{Remark}

\begin{Example}{Граничное поведение}{}
	Ряд $\sum_{n=1}^{\infty} \dfrac{z^n}{n}$ имеет $R = 1$ ($|a_n|^{1/n} = n^{-1/n} \to 1$): при $z = 1$ он расходится (гармонический ряд), а при $z = -1$ сходится (ряд Лейбница). Ряд $\sum_{n=1}^{\infty} \dfrac{z^n}{n^2}$ тоже имеет $R = 1$, но сходится абсолютно во всех точках окружности $|z| = 1$, так как мажорируется сходящимся рядом $\sum 1/n^2$.
\end{Example}
